%% RADiCAL Paper 2 (energy + position + 4D) -- NIM A (elsarticle) two-column style.
%% Compile: tectonic radical_energy_position.tex
%% STATUS: SKELETON. Section bodies are stubs; the localization analysis and final
%% figures still need to be built (see README.md and docs/PAPER2_DRAFT.md for the prose).
%% NOTE: the author block is a PLACEHOLDER; complete with the full RADiCAL author list
%% (Ref.~[2], NIM A 1068 (2024) 169737) before submission.
\documentclass[final,5p,times,twocolumn]{elsarticle}

\usepackage{graphicx}
\usepackage{amsmath}
\usepackage{booktabs}
\usepackage{xcolor}
\graphicspath{{figs/}}

\journal{Nuclear Instruments and Methods in Physics Research A}

\begin{document}

\begin{frontmatter}

\title{Shower-maximum energy and position measurement with wavelength-shifting
capillaries, and a simultaneous time--energy--position demonstration, in a RADiCAL
calorimeter module}

%% ---- Author block PLACEHOLDER (collaboration sets the full author order; see Ref.~[2]).
\author[uiowa]{J.~Wetzel\corref{cor}}
\ead{james@animal-lamps.com}
\cortext[cor]{Corresponding author.}
\author[radical]{\textit{for the RADiCAL Collaboration}}
\affiliation[uiowa]{organization={University of Iowa},
            city={Iowa City}, state={IA}, country={USA}}
\affiliation[radical]{organization={RADiCAL Collaboration}}
%% ---------------------------------------------------------------------------

\begin{abstract}
%% TODO. Same W/LYSO:Ce shashlik module and 25--150 GeV CERN SPS dataset as Paper~1.
%% Energy and transverse position from the four shower-maximum corner capillaries; the
%% E-type (full-length) energy capillary characterized in-beam; a simultaneous
%% time+energy+position ("4D") demonstration in one TENERGY module. Key numbers to land:
%% sigma_E/E ~14% @150 GeV (WLS-species-independent), position residual ~1.5 mm vs the
%% wire chamber, 4D demo sigma_t/sigma_E/sigma_x in one event sample.
\end{abstract}

\begin{keyword}
sampling calorimeter \sep shower maximum \sep wavelength shifter \sep position resolution
\sep 4D calorimetry \sep RADiCAL
\end{keyword}

\end{frontmatter}

%% ===========================================================================
\section{Introduction}
\label{sec:intro}
%% TODO: motivation; relation to the parent paper (NIM A 1068 (2024) 169737) §7.1, §7.3/7.4;
%% energy+position from the same shower-max signals; the 4D concept.

\section{Module and capillary types}
\label{sec:module}
%% TODO: brief apparatus recap (full spec: experimental_setup.md). T-type vs E-type capillaries;
%% TENERGY = 3 DSB1 T-type + 1 E-type. Reuse the Paper 1 setup figure/table as needed.

\section{Energy at shower maximum}
\label{sec:energy}
%% TODO: per-build sigma_E/E vs E (configResolutionFull.C -> config_sigmaE_vs_E). ~14% @150 GeV,
%% WLS-species-independent. Leakage-limited localized measurement; consistency with 52%/sqrt(E).

\section{Transverse position from the corner capillaries}
\label{sec:position}
%% TODO: 4-corner light center-of-gravity vs wire chamber (showerLocalization.C, transverseMaps.C).
%% Residual ~1.5 mm (1.52 x, 1.44 y) @150 GeV; module-ID concept.

\section{The E-type (full-length) energy capillary}
\label{sec:etype}
%% TODO: E-type vs T-type response & linearity (etypeEnergy.C, etypeChar.C). ~4-5x dimmer at
%% shower max, linear in energy; why it is excluded from the corner timing estimator.

\section{Simultaneous time, energy and position}
\label{sec:fourd}
%% TODO: the 4D demo in one TENERGY module (fourDdemo.C). sigma_t ~57 ps (best-bin ~35),
%% sigma_E/E ~11.9%, sigma_x ~0.9 mm from one event sample.

\section{Discussion and outlook}
\label{sec:discuss}
%% TODO: depth as the bridge between timing floor (Paper 1) and position signal; depth-corrected
%% estimator; enhanced modular design / 3x3 array (published §7.3/7.4).

\section{Conclusions}
\label{sec:conclusions}
%% TODO.

\section*{Acknowledgements}
We thank the CERN SPS staff for the test-beam operation. %% TODO: funding/agency text.

\begin{thebibliography}{9}
\bibitem{radical} C.~Perez-Lara, J.~Wetzel, U.~Akgun, et al. (RADiCAL Collaboration),
Nucl. Instrum. Methods A \textbf{1068} (2024) 169737.
%% TODO: complete references.
\end{thebibliography}

\end{document}
